\documentclass[11pt]{article}

\usepackage[T1]{fontenc}
\usepackage[utf8]{inputenc}
\usepackage[margin=1in]{geometry}
\usepackage{amsmath,amssymb,amsthm,mathtools}
\usepackage[colorlinks=true,linkcolor=blue,citecolor=blue,urlcolor=blue]{hyperref}

\numberwithin{equation}{section}
\newtheorem{theorem}{Theorem}[section]
\newtheorem{proposition}[theorem]{Proposition}
\newtheorem{lemma}[theorem]{Lemma}
\newtheorem{corollary}[theorem]{Corollary}
\newtheorem{conjecture}[theorem]{Conjecture}
\newtheorem{problem}[theorem]{Problem}
\theoremstyle{definition}
\newtheorem{definition}[theorem]{Definition}
\newtheorem{example}[theorem]{Example}
\theoremstyle{remark}
\newtheorem{remark}[theorem]{Remark}

\DeclareMathOperator{\Var}{Var}
\DeclareMathOperator{\Prob}{P}
\DeclareMathOperator{\supp}{supp}
\newcommand{\E}{\mathbb E}
\newcommand{\N}{\mathbb N}
\newcommand{\Z}{\mathbb Z}
\newcommand{\1}{\mathbf 1}

\title{Zeta-Law Entropy, Gamma Curvature, and Tail Partition Applications}
\author{Pudim AI Project}
\date{}

\begin{document}
\maketitle

\begin{abstract}
We study the Riemann zeta function as the partition function of the probability law \(\rho_\beta(n)=n^{-\beta}/\zeta(\beta)\) on the positive integers.  This normalization turns zeta ratios into moments, logarithmic derivatives into energy cumulants, divisor identities into probabilistic decompositions, and zeta tails into reciprocal partition problems.  The first main result identifies the microscopic successor entropy of the zeta law with the limit and supremum of finite modular successor entropies; for prime moduli these modular shadows are explicit nonlinear functionals of Dirichlet \(L\)-values.  The second structural result is an Euler-score decomposition expressing \(-(\log\zeta)'(\beta)\) as a von Mangoldt-weighted divisibility average.  We then prove six applications: the Alzer--Kwong convexity and concavity pattern for \(1/\zeta\) on negative real intervals; positivity of a Nantomah expression involving \(\zeta(n+1),\zeta(n+2),\zeta(n+3)\); a generalized Holder inequality for \(\Gamma(s)\zeta(s)\) extending Sroysang's formulation; complete monotonicity of \((\log\Gamma(x)+\log\Gamma(1/x))''\); and exact inverse-tail floor formulas for \(\zeta_n(7)^{-1}\) and \(\zeta_n(8)^{-1}\), where \(\zeta_n(s)=\sum_{k=n}^{\infty}k^{-s}\).  The framework gives a compact route from zeta probability laws to entropy, curvature, Mellin-integral inequalities, and discrete reciprocal partition formulas.
\end{abstract}

\noindent\textbf{Keywords.} Riemann zeta function; zeta probability law; Gibbs measure; entropy; modular residues; Dirichlet \(L\)-functions; von Mangoldt function; reciprocal zeta; zeta tails; Gamma function; polygamma functions; complete monotonicity; Holder inequality.

\section{Motivation and main results}

Many classical identities for \(\zeta(s)\) become more transparent after a probabilistic normalization.  For \(\beta>1\), the weights \(n^{-\beta}\) define a Gibbs law whose partition function is precisely \(\zeta(\beta)\).  The energy is \(\log n\), the free energy is \(\log\zeta(\beta)\), and ratios such as \(\zeta(\beta+r)/\zeta(\beta)\) become moments.  This viewpoint is elementary, but it packages several unrelated-looking zeta inequalities into a common calculus.

The paper develops five connected layers.  The moment layer records basic zeta-law identities.  The modular entropy layer compares the successor map \(n\mapsto n+1\) with its finite residue shadows modulo \(q\).  The Euler-score layer rewrites zeta energy through von Mangoldt divisibility events.  The Mellin-Planck layer treats \(\Gamma(s)\zeta(s)\) as a continuous partition function and converts generalized Holder inequalities into integral norm inequalities.  The tail layer studies \(\zeta_n(s)^{-1}\) as a reciprocal partition problem for the truncated zeta law.

Section~\ref{sec:applications} records the six solved external problems or candidate-problem specializations with stable application labels.  APP-0001 is the Alzer--Kwong reciprocal-zeta conjecture, proved by transporting a positive-axis curvature estimate through the functional equation.  APP-0002 is the Nantomah positivity question, proved by splitting a zeta expression into a positive linear tail and a controlled quadratic correction.  APP-0003 is the Sroysang generalized Holder problem, proved directly from generalized Holder applied to the Planck kernel.  APP-0004 is a reciprocal-Gamma complete-monotonicity problem, solved by a Weierstrass-product Laplace kernel.  APP-0005 and APP-0006 give exact formulas for \(\lfloor \zeta_n(7)^{-1}\rfloor\) and \(\lfloor \zeta_n(8)^{-1}\rfloor\), resolving the first two local tail-inversion special cases while leaving the general integer \(s>6\) problem open.

\section{Applications}
\label{sec:applications}

The following application labels are used consistently in the repository landing page, the application changelog, and the theorem names below.  The order here is the landing-page order; the proof bodies are placed where their methods naturally fit the paper.

\begin{enumerate}
\item \textbf{APP-0001.} Alzer--Kwong convexity and concavity problem.  Solved by Theorem~\ref{thm:alzer-kwong}; source problem reference \cite{alzer-kwong}.
\item \textbf{APP-0002.} Nantomah zeta positivity problem.  Solved by Theorem~\ref{thm:nantomah-positivity}; source problem reference \cite{nantomah}.
\item \textbf{APP-0003.} Sroysang generalized Holder problem.  Solved by Theorem~\ref{thm:sroysang-holder}; source problem reference \cite{sroysang}.
\item \textbf{APP-0004.} Complete monotonicity of \((\log\Gamma(x)+\log\Gamma(1/x))''\).  Solved by Theorem~\ref{thm:reciprocal-gamma-complete-monotonicity}; source problem reference \cite{qi-lim-nantomah-polygamma}.
\item \textbf{APP-0005.} Exact inverse-tail floor formula at \(s=7\).  Solved by Theorem~\ref{thm:tail-s7-floor}; source context \cite{kim-song-tail,pan-wu-tail}.
\item \textbf{APP-0006.} Exact inverse-tail floor formula at \(s=8\).  Solved by Theorem~\ref{thm:tail-s8-floor}; source context \cite{kim-song-tail,pan-wu-tail}.
\end{enumerate}

\section{Notation and zeta-law calculus}

\begin{definition}[Riemann zeta probability law]
\label{def:zeta-law}
For \(\beta>1\), define a probability law \(\rho_\beta\) on \(\N\) by
\begin{equation}
\label{eq:zeta-law}
\rho_\beta(n)=\frac{n^{-\beta}}{\zeta(\beta)},\qquad n\in\N.
\end{equation}
Equivalently, with \(E(n)=\log n\) and \(Z(\beta)=\zeta(\beta)\),
\begin{equation}
\label{eq:gibbs-form}
\rho_\beta(n)=\frac{e^{-\beta E(n)}}{Z(\beta)}.
\end{equation}
\end{definition}

\begin{definition}[Zeta free energy]
\label{def:zeta-free-energy}
The zeta free energy is
\begin{equation}
\label{eq:zeta-free-energy}
A(\beta)=\log\zeta(\beta),\qquad \beta>1.
\end{equation}
\end{definition}

\begin{proposition}[Free-energy derivatives]
\label{prop:free-energy-derivatives}
For every \(\beta>1\),
\begin{align}
\label{eq:free-energy-derivatives}
A'(\beta)&=-\E_\beta[\log N],&
A''(\beta)&=\Var_\beta(\log N).
\end{align}
\end{proposition}

\begin{proof}
Absolute convergence of the differentiated Dirichlet series on compact subintervals of \((1,\infty)\) permits termwise differentiation.  Thus
\begin{equation}
\label{eq:first-free-derivative-proof}
A'(\beta)=\frac{\zeta'(\beta)}{\zeta(\beta)}
=-\frac{\sum_{n\ge1}(\log n)n^{-\beta}}{\zeta(\beta)}
=-\E_\beta[\log N].
\end{equation}
Differentiating once more gives
\begin{equation}
\label{eq:second-free-derivative-proof}
A''(\beta)=\E_\beta[(\log N)^2]-\E_\beta[\log N]^2
=\Var_\beta(\log N).
\end{equation}
\end{proof}

\begin{proposition}[Zeta-law calculus]
\label{prop:zeta-law-calculus}
Let \(\beta>1\), and let \(N\) have law \(\rho_\beta\).  For every \(r\ge0\),
\begin{equation}
\label{eq:zeta-law-moment}
\E_\beta[N^{-r}]=\frac{\zeta(\beta+r)}{\zeta(\beta)}.
\end{equation}
Moreover, for \(\alpha,\beta>1\),
\begin{equation}
\label{eq:zeta-law-relative-entropy}
D(\rho_\alpha\Vert\rho_\beta)
=(\beta-\alpha)\E_\alpha[\log N]
+\log\zeta(\beta)-\log\zeta(\alpha).
\end{equation}
\end{proposition}

\begin{proof}
The moment identity follows by direct summation:
\begin{equation}
\label{eq:zeta-law-moment-proof}
\E_\beta[N^{-r}]
=\sum_{n=1}^{\infty}n^{-r}\frac{n^{-\beta}}{\zeta(\beta)}
=\frac{\zeta(\beta+r)}{\zeta(\beta)}.
\end{equation}
For relative entropy, substitute \eqref{eq:zeta-law}:
\begin{equation}
\label{eq:zeta-law-relative-entropy-proof}
D(\rho_\alpha\Vert\rho_\beta)
=\sum_{n=1}^{\infty}\rho_\alpha(n)
\left((\beta-\alpha)\log n+\log\zeta(\beta)-\log\zeta(\alpha)\right),
\end{equation}
which is \eqref{eq:zeta-law-relative-entropy}.
\end{proof}

\section{Tail zeta laws and inverse-tail floors}

\begin{definition}[Tail zeta partition function]
\label{def:tail-zeta-partition}
For \(s>1\) and \(n\in\N\), define the tail zeta partition function by
\begin{equation}
\label{eq:tail-zeta-partition}
Z_{s,n}=\zeta_n(s)=\sum_{k=n}^{\infty}k^{-s}.
\end{equation}
The corresponding tail zeta law is the probability distribution on \(\{n,n+1,\ldots\}\) given by
\begin{equation}
\label{eq:tail-zeta-law}
\rho_{s,n}(k)=\frac{k^{-s}}{\zeta_n(s)},\qquad k\ge n.
\end{equation}
\end{definition}

Exact formulas for integer parts of reciprocal zeta tails have been studied for small integer exponents.  Known formulas include \(s=2,3,4,5\), with a sufficiently-large-\(n\) formula for \(s=6\); Kim and Song record that no analogous formula was known for integer \(s>6\) at the time of their work \cite{kim-song-tail}.  The 2024 asymptotic inverse-tail results of Pan and Wu provide related context \cite{pan-wu-tail}.  The next two theorems give exact formulas at \(s=7\) and \(s=8\).  The first uses a single rational telescoping enclosure, while the second needs adjacent corrected asymptotic truncations.

\begin{theorem}[APP-0005: Exact inverse-tail floor formula at s=7]
\label{thm:tail-s7-floor}
Let
\begin{equation}
\label{eq:tail-s7}
T_7(n)=\zeta_n(7)=\sum_{k=n}^{\infty}\frac1{k^7}.
\end{equation}
Define
\begin{equation}
\label{eq:tail-s7-q}
Q(n)=120n^6-360n^5+660n^4-720n^3+354n^2-54n+375,
\qquad
P(n)=\frac{Q(n)}{20}.
\end{equation}
Then, for every \(n\ge28\),
\begin{equation}
\label{eq:tail-s7-floor}
\left\lfloor T_7(n)^{-1}\right\rfloor=\left\lfloor P(n)\right\rfloor.
\end{equation}
For \(1\le n\le27\), the exact values are
\begin{equation}
\label{eq:tail-s7-finite-table}
\begin{array}{c|r@{\qquad}c|r@{\qquad}c|r}
n&\lfloor T_7(n)^{-1}\rfloor&n&\lfloor T_7(n)^{-1}\rfloor&n&\lfloor T_7(n)^{-1}\rfloor\\
\hline
1&0&10&4495760&19&241765529\\
2&119&11&8167814&20&331399044\\
3&1862&12&14061542&21&447175152\\
4&12573&13&23143975&22&594869693\\
5&54069&14&36668777&23&781167220\\
6&175597&15&56228085&24&1013751716\\
7&471118&16&83808666&25&1301401638\\
8&1100904&17&121852400&26&1654089273\\
9&2317459&18&173321080&27&2083084422
\end{array}.
\end{equation}
\end{theorem}

\begin{proof}
The polynomial \(Q\) is positive on \(\N\), because
\begin{equation}
\label{eq:tail-s7-q-positive}
\begin{aligned}
Q(n)
&=120(n-1)^6+360(n-1)^5+660(n-1)^4+720(n-1)^3\\
&\qquad +354(n-1)^2+54(n-1)+375.
\end{aligned}
\end{equation}
For \(k\ge1\), direct expansion gives
\begin{equation}
\label{eq:tail-s7-lower-difference}
\frac1{k^7}
-\left(\frac{20}{Q(k)}-\frac{20}{Q(k+1)}\right)
=
\frac{9(60284k^4+29176k^2+15625)}
{k^7Q(k)Q(k+1)}>0.
\end{equation}
Summing \eqref{eq:tail-s7-lower-difference} over \(k\ge n\) yields
\begin{equation}
\label{eq:tail-s7-lower-sum}
T_7(n)>\frac{20}{Q(n)}=\frac1{P(n)},
\end{equation}
and therefore \(T_7(n)^{-1}<P(n)\).

For \(k\ge28\), another expansion gives
\begin{equation}
\label{eq:tail-s7-upper-difference}
\left(\frac{20}{Q(k)-3}-\frac{20}{Q(k+1)-3}\right)
-\frac1{k^7}
=
\frac{36(20k^6-14961k^4-7235k^2-3844)}
{k^7(Q(k)-3)(Q(k+1)-3)}.
\end{equation}
The numerator in \eqref{eq:tail-s7-upper-difference} is positive for \(k\ge28\).  Indeed, writing \(k=m+28\), it becomes
\begin{equation}
\label{eq:tail-s7-upper-positive}
20m^6+3360m^5+220239m^4+7105168m^3
+114013021m^2+751143512m+436261580,
\end{equation}
which has positive coefficients for \(m\ge0\).  Hence, for \(n\ge28\),
\begin{equation}
\label{eq:tail-s7-upper-sum}
T_7(n)<\frac{20}{Q(n)-3}=\frac1{P(n)-3/20}.
\end{equation}
Combining \eqref{eq:tail-s7-lower-sum} and \eqref{eq:tail-s7-upper-sum} gives
\begin{equation}
\label{eq:tail-s7-inverse-window}
P(n)-\frac3{20}<T_7(n)^{-1}<P(n).
\end{equation}

Modulo \(20\),
\begin{equation}
\label{eq:tail-s7-q-mod}
Q(n)\equiv 14n^2-14n+15\pmod {20}.
\end{equation}
Since \(n(n-1)\pmod {10}\in\{0,2,6\}\), this implies
\begin{equation}
\label{eq:tail-s7-q-residues}
Q(n)\pmod {20}\in\{3,15,19\}.
\end{equation}
Thus the fractional part of \(P(n)\) is one of \(3/20,15/20,19/20\).  If \(a=\lfloor P(n)\rfloor\), then \(P(n)-3/20\ge a\), and \eqref{eq:tail-s7-inverse-window} gives
\begin{equation}
\label{eq:tail-s7-floor-trap}
a<T_7(n)^{-1}<a+1.
\end{equation}
This proves \eqref{eq:tail-s7-floor} for \(n\ge28\).

For the finite range, \(T_7(1)>1\), so \(0<T_7(1)^{-1}<1\).  For \(2\le n\le27\), let \(M_n\) denote the corresponding value in \eqref{eq:tail-s7-finite-table}.  With \(N=1500\), exact rational arithmetic verifies
\begin{equation}
\label{eq:tail-s7-finite-lower}
\sum_{k=n}^{N}\frac1{k^7}>\frac1{M_n+1}
\end{equation}
and
\begin{equation}
\label{eq:tail-s7-finite-upper}
\sum_{k=n}^{N}\frac1{k^7}+\frac1{6N^6}<\frac1{M_n}.
\end{equation}
The monotone integral estimate
\begin{equation}
\label{eq:tail-s7-integral-tail}
\sum_{k=N+1}^{\infty}\frac1{k^7}<\int_N^\infty x^{-7}\,dx=\frac1{6N^6}
\end{equation}
then gives
\begin{equation}
\label{eq:tail-s7-finite-window}
\frac1{M_n+1}<T_7(n)<\frac1{M_n}.
\end{equation}
Therefore \(M_n<T_7(n)^{-1}<M_n+1\), proving \eqref{eq:tail-s7-finite-table}.  The finite rational certificate is stored with the source as a reproducible calculation.
\end{proof}

\begin{theorem}[APP-0006: Exact inverse-tail floor formula at s=8]
\label{thm:tail-s8-floor}
Let
\begin{equation}
\label{eq:tail-s8}
T_8(n)=\zeta_n(8)=\sum_{k=n}^{\infty}\frac1{k^8}.
\end{equation}
Define
\begin{align}
\label{eq:tail-s8-a}
A_8(n)
={}&7n^7-\frac{49}{2}n^6+\frac{637}{12}n^5-\frac{1715}{24}n^4
+\frac{7399}{144}n^3-\frac{1715}{96}n^2+\frac{69685}{1728}n-\frac{65611}{3456}\notag\\
&-\frac{23764069}{103680n}
-\frac{23764069}{207360n^2}
+\frac{2102131409}{1244160n^3}
+\frac{2149659547}{829440n^4}\notag\\
&-\frac{203290800013}{14929920n^5}
-\frac{1146003910541}{29859840n^6}
+\frac{21587525169497}{179159040n^7}.
\end{align}
Then, for every \(n\ge6\),
\begin{equation}
\label{eq:tail-s8-floor}
\left\lfloor T_8(n)^{-1}\right\rfloor=\left\lfloor A_8(n)\right\rfloor.
\end{equation}
For \(1\le n\le5\), the exact values are
\begin{equation}
\label{eq:tail-s8-finite-table}
\begin{array}{c|rrrrr}
n&1&2&3&4&5\\
\hline
\lfloor T_8(n)^{-1}\rfloor&0&245&5844&53503&291407
\end{array}.
\end{equation}
\end{theorem}

\begin{proof}
Set
\begin{equation}
\label{eq:tail-s8-b}
B_8(n)=A_8(n)+\frac{200335112892761}{358318080n^8}.
\end{equation}
For \(C=A_8\) or \(B_8\), write \(R_C(k)=1/C(k)\).  Exact denominator clearing gives the two telescoping inequalities
\begin{equation}
\label{eq:tail-s8-lower-inverse-difference}
R_{A_8}(k)-R_{A_8}(k+1)-\frac1{k^8}>0\qquad(k\ge3)
\end{equation}
and
\begin{equation}
\label{eq:tail-s8-upper-inverse-difference}
\frac1{k^8}-\left(R_{B_8}(k)-R_{B_8}(k+1)\right)>0\qquad(k\ge1).
\end{equation}
In \eqref{eq:tail-s8-lower-inverse-difference}, the cleared numerator is a polynomial in \(k-3\) with all coefficients positive; in \eqref{eq:tail-s8-upper-inverse-difference}, the cleared numerator is a polynomial in \(k-1\) with all coefficients positive.  The exact coefficient verification is recorded in the rational certificate accompanying the manuscript.

Summing \eqref{eq:tail-s8-lower-inverse-difference} over \(k\ge n\) gives
\begin{equation}
\label{eq:tail-s8-lower-inverse-sum}
T_8(n)<R_{A_8}(n)=\frac1{A_8(n)}
\end{equation}
for \(n\ge3\), and hence \(A_8(n)<T_8(n)^{-1}\).  Summing \eqref{eq:tail-s8-upper-inverse-difference} gives
\begin{equation}
\label{eq:tail-s8-upper-inverse-sum}
T_8(n)>R_{B_8}(n)=\frac1{B_8(n)},
\end{equation}
and hence \(T_8(n)^{-1}<B_8(n)\).  Therefore, for \(n\ge6\),
\begin{equation}
\label{eq:tail-s8-inverse-window}
A_8(n)<T_8(n)^{-1}<B_8(n).
\end{equation}

It remains to see that no integer lies between \(A_8(n)\) and \(B_8(n)\).  Exact integer arithmetic verifies
\begin{equation}
\label{eq:tail-s8-floor-gap-finite}
\lfloor A_8(n)\rfloor=\lfloor B_8(n)\rfloor\qquad(6\le n\le10^6).
\end{equation}
For \(n>10^6\), let \(P_8(n)\) be the polynomial part of \(A_8(n)\), through the constant term.  The numerator of \(P_8(n)\) modulo \(3456\) runs through the odd residue classes, so \(P_8(n)\) is at distance at least \(1/3456\) from the next integer.  The correction \(A_8(n)-P_8(n)\) is negative and has absolute value \(<1/3456\), while
\begin{equation}
\label{eq:tail-s8-gap-small}
0<B_8(n)-A_8(n)<\frac1{3456}.
\end{equation}
Thus \(B_8(n)\) is still below the next integer after \(A_8(n)\).  Combining this with \eqref{eq:tail-s8-inverse-window} proves \eqref{eq:tail-s8-floor}.

For \(1\le n\le5\), \eqref{eq:tail-s8-finite-table} is certified by exact rational partial sums.  With \(N=300\), the displayed value \(M_n\) satisfies, for \(2\le n\le5\),
\begin{equation}
\label{eq:tail-s8-finite-lower}
\sum_{k=n}^{N}\frac1{k^8}>\frac1{M_n+1}
\end{equation}
and
\begin{equation}
\label{eq:tail-s8-finite-upper}
\sum_{k=n}^{N}\frac1{k^8}+\frac1{7N^7}<\frac1{M_n}.
\end{equation}
Since
\begin{equation}
\label{eq:tail-s8-integral-tail}
\sum_{k=N+1}^{\infty}\frac1{k^8}<\int_N^\infty x^{-8}\,dx=\frac1{7N^7},
\end{equation}
we get \(M_n<T_8(n)^{-1}<M_n+1\).  For \(n=1\), \(T_8(1)>1\), so the floor is \(0\).
\end{proof}

\section{Modular successor entropy}

\begin{definition}[Modular residue distribution and successor entropy]
\label{def:modular-residue-successor}
For \(q\ge2\), define the residue distribution of \(\rho_\beta\) modulo \(q\) by
\begin{equation}
\label{eq:modular-residue-distribution}
\mu_{q,\beta}(a)=
\sum_{\substack{n\ge1\\ n\equiv a\pmod q}}\rho_\beta(n),
\qquad a\in\Z/q\Z.
\end{equation}
Writing \(a+1\) cyclically in \(\Z/q\Z\), define
\begin{equation}
\label{eq:modular-successor-entropy}
B_q(\beta)=
\sum_{a\in\Z/q\Z}
\mu_{q,\beta}(a)\log\frac{\mu_{q,\beta}(a)}{\mu_{q,\beta}(a+1)}.
\end{equation}
\end{definition}

\begin{theorem}[Zeta-law successor entropy and modular resolution]
\label{thm:successor-entropy}
For every \(\beta>1\), define
\begin{equation}
\label{eq:successor-entropy}
\Sigma(\beta)=\sum_{n=1}^{\infty}\rho_\beta(n)
\log\frac{\rho_\beta(n)}{\rho_\beta(n+1)}.
\end{equation}
Then
\begin{equation}
\label{eq:successor-entropy-series}
\Sigma(\beta)=
\frac{\beta}{\zeta(\beta)}
\sum_{k=1}^{\infty}\frac{(-1)^{k+1}}{k}\zeta(\beta+k),
\end{equation}
and
\begin{equation}
\label{eq:successor-entropy-modular-resolution}
\Sigma(\beta)=\sup_{q\ge2}B_q(\beta)=\lim_{q\to\infty}B_q(\beta).
\end{equation}
\end{theorem}

\begin{proof}
Since
\begin{equation}
\label{eq:successor-ratio}
\frac{\rho_\beta(n)}{\rho_\beta(n+1)}
=\left(1+\frac1n\right)^\beta,
\end{equation}
we have
\begin{equation}
\label{eq:successor-entropy-log-sum}
\Sigma(\beta)=\frac{\beta}{\zeta(\beta)}
\sum_{n=1}^{\infty}n^{-\beta}\log\left(1+\frac1n\right).
\end{equation}
Using
\begin{equation}
\label{eq:log-series}
\log(1+x)=\sum_{k=1}^{\infty}\frac{(-1)^{k+1}}{k}x^k
\qquad(0<x\le1)
\end{equation}
with Abel limiting at \(x=1\), and summing against \(n^{-\beta}\), gives \eqref{eq:successor-entropy-series}.

For the modular identity, fix \(q\ge2\).  The log-sum inequality gives, for each residue class \(a\),
\begin{equation}
\label{eq:modular-log-sum}
\sum_{\substack{n\ge1\\ n\equiv a\pmod q}}
\rho_\beta(n)\log\frac{\rho_\beta(n)}{\rho_\beta(n+1)}
\ge
\mu_{q,\beta}(a)
\log\frac{\mu_{q,\beta}(a)}
{\sum_{n\equiv a\pmod q}\rho_\beta(n+1)}.
\end{equation}
For \(a\ne0\), the denominator is \(\mu_{q,\beta}(a+1)\).  For \(a=0\), it is \(\mu_{q,\beta}(1)-\rho_\beta(1)\le\mu_{q,\beta}(1)\).  Summing \eqref{eq:modular-log-sum} over \(a\) gives
\begin{equation}
\label{eq:successor-dominates-modular}
\Sigma(\beta)\ge B_q(\beta)
\qquad(q\ge2).
\end{equation}

Conversely, fix \(M\in\N\).  If \(q>M+1\), then for \(1\le a\le M+1\),
\begin{equation}
\label{eq:residue-isolation}
\mu_{q,\beta}(a)=\rho_\beta(a)+O_{\beta,M}(q^{-\beta}),
\end{equation}
and \(\mu_{q,\beta}(0)=q^{-\beta}\).  All cyclic terms in \(B_q(\beta)\) are nonnegative except possibly the boundary term, which satisfies
\begin{equation}
\label{eq:boundary-term}
\mu_{q,\beta}(0)\log\frac{\mu_{q,\beta}(0)}{\mu_{q,\beta}(1)}
=O_\beta(q^{-\beta}\log q).
\end{equation}
Therefore
\begin{equation}
\label{eq:modular-liminf}
\liminf_{q\to\infty}B_q(\beta)
\ge
\sum_{n=1}^{M}\rho_\beta(n)\log\frac{\rho_\beta(n)}{\rho_\beta(n+1)}.
\end{equation}
Letting \(M\to\infty\) gives \(\liminf_qB_q(\beta)\ge\Sigma(\beta)\), which together with \eqref{eq:successor-dominates-modular} proves \eqref{eq:successor-entropy-modular-resolution}.
\end{proof}

\begin{corollary}[Prime-modulus Dirichlet L-resolution]
\label{cor:prime-dirichlet-resolution}
Let \(p\) be prime, let \(\chi\) range over Dirichlet characters modulo \(p\), and define
\begin{equation}
\label{eq:prime-dirichlet-shadow}
S_a(p,\beta)=\sum_{\chi\bmod p}\overline{\chi(a)}L(\beta,\chi)
\qquad(1\le a\le p-1).
\end{equation}
Then
\begin{equation}
\label{eq:prime-residue-formula}
\mu_{p,\beta}(a)=\frac{S_a(p,\beta)}{(p-1)\zeta(\beta)}
\quad(1\le a\le p-1),
\qquad
\mu_{p,\beta}(0)=p^{-\beta}.
\end{equation}
Consequently the prime-modulus functionals obtained by substituting \eqref{eq:prime-residue-formula} into \(B_p(\beta)\) satisfy
\begin{equation}
\label{eq:prime-functional-limit}
\Sigma(\beta)=
\lim_{\substack{p\to\infty\\ p\ {\rm prime}}}
F_p\bigl(\zeta(\beta),\{L(\beta,\chi)\}_{\chi\bmod p}\bigr)
=
\sup_{p\ {\rm prime}}
F_p\bigl(\zeta(\beta),\{L(\beta,\chi)\}_{\chi\bmod p}\bigr).
\end{equation}
\end{corollary}

\begin{proof}
For \(a\not\equiv0\pmod p\), character orthogonality gives
\begin{equation}
\label{eq:character-orthogonality}
\frac1{p-1}\sum_{\chi\bmod p}\overline{\chi(a)}\chi(n)
=
\begin{cases}
1,& n\equiv a\pmod p,\\
0,& n\not\equiv a\pmod p.
\end{cases}
\end{equation}
Therefore
\begin{equation}
\label{eq:dirichlet-residue-sum}
\sum_{\substack{n\ge1\\ n\equiv a\pmod p}}n^{-\beta}
=\frac1{p-1}\sum_{\chi\bmod p}\overline{\chi(a)}L(\beta,\chi).
\end{equation}
Dividing by \(\zeta(\beta)\) gives the nonzero residue formula, while
\begin{equation}
\label{eq:zero-residue-prime}
\mu_{p,\beta}(0)=\frac1{\zeta(\beta)}\sum_{m=1}^{\infty}(mp)^{-\beta}
=p^{-\beta}.
\end{equation}
The functional identity follows by substituting these residue formulas into \(B_p(\beta)\) and applying Theorem \ref{thm:successor-entropy} along the prime sequence.
\end{proof}

\section{Euler-score identities}

\begin{theorem}[Euler-score decomposition]
\label{thm:euler-score-decomposition}
For every \(\beta>1\),
\begin{equation}
\label{eq:euler-score-decomposition}
-\bigl(\log\zeta\bigr)'(\beta)
=\E_\beta[\log N]
=\sum_{d=1}^{\infty}\Lambda(d)\mu_{d,\beta}(0)
=\sum_{d=1}^{\infty}\frac{\Lambda(d)}{d^\beta}.
\end{equation}
Equivalently,
\begin{equation}
\label{eq:log-zeta-von-mangoldt}
-\frac{\zeta'(\beta)}{\zeta(\beta)}
=\sum_{d=1}^{\infty}\frac{\Lambda(d)}{d^\beta}.
\end{equation}
\end{theorem}

\begin{proof}
The von Mangoldt identity is
\begin{equation}
\label{eq:mangoldt-identity}
\log n=\sum_{d\mid n}\Lambda(d).
\end{equation}
By nonnegative summation,
\begin{equation}
\label{eq:euler-score-proof}
\E_\beta[\log N]
=\sum_{n\ge1}\rho_\beta(n)\sum_{d\mid n}\Lambda(d)
=\sum_{d\ge1}\Lambda(d)\Prob_\beta(d\mid N).
\end{equation}
The divisibility probability is
\begin{equation}
\label{eq:divisibility-probability}
\Prob_\beta(d\mid N)
=\frac1{\zeta(\beta)}\sum_{m\ge1}(dm)^{-\beta}
=d^{-\beta}.
\end{equation}
Substitution into \eqref{eq:euler-score-proof} proves \eqref{eq:euler-score-decomposition}; \eqref{eq:log-zeta-von-mangoldt} follows from \eqref{eq:first-free-derivative-proof}.
\end{proof}

\begin{proposition}[Finite Euler-score identity]
\label{prop:finite-euler-score}
For every \(A\subseteq\{1,\ldots,N\}\),
\begin{equation}
\label{eq:finite-euler-score}
\sum_{n\in A}\log n
=\sum_{d\le N}\Lambda(d)\sum_{m\le N/d}\1_A(md).
\end{equation}
\end{proposition}

\begin{proof}
Reverse the order of summation:
\begin{equation}
\label{eq:finite-euler-score-proof}
\sum_{d\le N}\Lambda(d)\sum_{m\le N/d}\1_A(md)
=\sum_{n\in A}\sum_{d\mid n}\Lambda(d)
=\sum_{n\in A}\log n.
\end{equation}
\end{proof}

\section{Curvature and Mellin-Planck tools}

\begin{lemma}[Uniform positive-axis curvature bound]
\label{lem:uniform-positive-axis-curvature}
For every \(s\ge3\),
\begin{equation}
\label{eq:uniform-curvature-bound}
(\log\zeta)''(s)=\Var_s(\log N)<\frac13.
\end{equation}
\end{lemma}

\begin{proof}
By Proposition \ref{prop:free-energy-derivatives},
\begin{equation}
\label{eq:curvature-first-bound}
(\log\zeta)''(s)
=\frac{\zeta''(s)}{\zeta(s)}
-\left(\frac{\zeta'(s)}{\zeta(s)}\right)^2
\le \frac{\zeta''(s)}{\zeta(s)}
\le \zeta''(s).
\end{equation}
For \(s\ge3\),
\begin{equation}
\label{eq:zeta-second-derivative-bound}
\zeta''(s)=\sum_{m=2}^{\infty}\frac{(\log m)^2}{m^s}
\le\sum_{m=2}^{\infty}\frac{(\log m)^2}{m^3}.
\end{equation}
The function \(x\mapsto(\log x)^2x^{-3}\) is decreasing for \(x\ge2\), so
\begin{equation}
\label{eq:curvature-integral-comparison}
\sum_{m=2}^{\infty}\frac{(\log m)^2}{m^3}
\le
\frac{(\log2)^2}{8}
+\int_2^\infty\frac{(\log x)^2}{x^3}\,dx.
\end{equation}
Since
\begin{equation}
\label{eq:curvature-integral-value}
\int_2^\infty\frac{(\log x)^2}{x^3}\,dx
=\frac{(\log2)^2+\log2+\frac12}{8},
\end{equation}
we obtain
\begin{equation}
\label{eq:curvature-final-bound}
(\log\zeta)''(s)
\le
\frac{2(\log2)^2+\log2+\frac12}{8}
<\frac13.
\end{equation}
\end{proof}

\begin{definition}[Mellin-Planck partition function]
\label{def:mellin-planck-partition}
For \(s>1\), define
\begin{equation}
\label{eq:mellin-planck-partition}
M(s)=\Gamma(s)\zeta(s)
=\int_0^\infty\frac{t^{s-1}}{e^t-1}\,dt.
\end{equation}
The corresponding probability law is
\begin{equation}
\label{eq:mellin-planck-law}
\nu_s(dt)=\frac{t^{s-1}}{(e^t-1)M(s)}\,dt.
\end{equation}
\end{definition}

\begin{proposition}[Mellin-Planck log-convexity]
\label{prop:mellin-planck-log-convexity}
For \(s>1\),
\begin{equation}
\label{eq:mellin-planck-log-convexity}
\frac{d^2}{ds^2}\log M(s)=\Var_{\nu_s}(\log t)\ge0.
\end{equation}
\end{proposition}

\begin{proof}
Differentiate the logarithm of the integral in \eqref{eq:mellin-planck-partition} under the integral sign.  The first derivative is the \(\nu_s\)-expectation of \(\log t\), and the second derivative is its variance:
\begin{equation}
\label{eq:mellin-planck-variance-proof}
\frac{d^2}{ds^2}\log M(s)
=\E_{\nu_s}[(\log t)^2]-\E_{\nu_s}[\log t]^2.
\end{equation}
\end{proof}

Qi, Lim, and Nantomah asked whether the second derivative of \(\log\Gamma(x)+\log\Gamma(1/x)\) is completely monotonic on \((0,\infty)\) \cite{qi-lim-nantomah-polygamma}.  The next theorem answers this in the affirmative by a reciprocal-Weierstrass Laplace kernel.

\begin{theorem}[APP-0004: Complete monotonicity of reciprocal-Gamma curvature]
\label{thm:reciprocal-gamma-complete-monotonicity}
Let
\begin{equation}
\label{eq:reciprocal-gamma-h}
H(x)=\log\Gamma(x)+\log\Gamma(1/x),\qquad x>0.
\end{equation}
Then \(H''\) is strictly completely monotonic on \((0,\infty)\); that is,
\begin{equation}
\label{eq:reciprocal-gamma-complete-monotonicity}
(-1)^m\frac{d^m}{dx^m}H''(x)>0
\qquad(m\ge0,\ x>0).
\end{equation}
\end{theorem}

\begin{proof}
The Weierstrass product gives, for \(x>0\),
\begin{equation}
\label{eq:gamma-weierstrass-log}
\log\Gamma(x)
=-\log x-\gamma x-\sum_{n=1}^{\infty}
\left(\log\left(1+\frac{x}{n}\right)-\frac{x}{n}\right).
\end{equation}
After substituting \(1/x\) in \eqref{eq:gamma-weierstrass-log}, adding the two formulas, and differentiating twice on compact subintervals of \((0,\infty)\), the logarithmic second-derivative terms cancel and
\begin{equation}
\label{eq:reciprocal-gamma-second-series}
\begin{aligned}
H''(x)
&=\sum_{n=1}^{\infty}\frac1{(x+n)^2}
-\frac{2\gamma}{x^3}\\
&\qquad+
\sum_{n=1}^{\infty}
\left[
\frac1{(x+1/n)^2}
-\frac1{x^2}
+\frac{2}{nx^3}
\right].
\end{aligned}
\end{equation}
The elementary Laplace identities
\begin{equation}
\label{eq:reciprocal-gamma-laplace-identities}
\frac1{(x+a)^2}=\int_0^\infty t e^{-(x+a)t}\,dt,\qquad
\frac1{x^2}=\int_0^\infty t e^{-xt}\,dt,\qquad
\frac{2a}{x^3}=a\int_0^\infty t^2e^{-xt}\,dt
\end{equation}
show that, with \(a=1/n\),
\begin{equation}
\label{eq:reciprocal-gamma-defect-laplace}
\frac1{(x+1/n)^2}-\frac1{x^2}+\frac{2}{nx^3}
=\int_0^\infty e^{-xt}t\left(e^{-t/n}-1+\frac{t}{n}\right)\,dt.
\end{equation}
Also,
\begin{equation}
\label{eq:reciprocal-gamma-trigamma-laplace}
\sum_{n=1}^{\infty}\frac1{(x+n)^2}
=\int_0^\infty e^{-xt}\frac{t}{e^t-1}\,dt,
\end{equation}
and
\begin{equation}
\label{eq:reciprocal-gamma-euler-laplace}
-\frac{2\gamma}{x^3}
=\int_0^\infty e^{-xt}(-\gamma t^2)\,dt.
\end{equation}
Combining \eqref{eq:reciprocal-gamma-second-series}--\eqref{eq:reciprocal-gamma-euler-laplace} gives
\begin{equation}
\label{eq:reciprocal-gamma-kernel-representation}
H''(x)=\int_0^\infty e^{-xt}K(t)\,dt,
\end{equation}
where
\begin{equation}
\label{eq:reciprocal-gamma-kernel}
K(t)=\frac{t}{e^t-1}
+t\sum_{n=1}^{\infty}\left(e^{-t/n}-1+\frac{t}{n}\right)
-\gamma t^2.
\end{equation}

It remains to prove \(K(t)>0\).  Put
\begin{equation}
\label{eq:reciprocal-gamma-phi}
\phi(u)=e^{-u}-1+u.
\end{equation}
Then \(\phi\ge0\), \(\phi'=1-e^{-u}\ge0\), and \(u\mapsto\phi(t/u)\) is positive and decreasing on \([1,\infty)\).  Therefore
\begin{equation}
\label{eq:reciprocal-gamma-sum-integral}
\sum_{n=1}^{\infty}\phi(t/n)\ge\int_1^\infty \phi(t/u)\,du.
\end{equation}
With \(v=t/u\), integration by parts gives
\begin{equation}
\label{eq:reciprocal-gamma-integral-transform}
\int_1^\infty\phi(t/u)\,du
=t\int_0^t\frac{\phi(v)}{v^2}\,dv
=-\phi(t)+t\int_0^t\frac{1-e^{-v}}{v}\,dv.
\end{equation}
Using
\begin{equation}
\label{eq:reciprocal-gamma-e1-identity}
\int_0^t\frac{1-e^{-v}}{v}\,dv
=\gamma+\log t+E_1(t),
\qquad
E_1(t)=\int_t^\infty\frac{e^{-v}}{v}\,dv,
\end{equation}
we obtain
\begin{equation}
\label{eq:reciprocal-gamma-kernel-lower-bound}
\frac{K(t)}{t}
\ge
\frac1{e^t-1}-\phi(t)+t\log t+tE_1(t).
\end{equation}
Since \(E_1(t)>0\) and \(\log t\ge1-1/t\) for \(t>0\),
\begin{equation}
\label{eq:reciprocal-gamma-kernel-positive}
\frac{K(t)}{t}
>
\frac1{e^t-1}-\phi(t)+t-1
=\frac1{e^t-1}-e^{-t}
=\frac{e^{-t}}{e^t-1}>0.
\end{equation}
Thus \(K(t)>0\) for \(t>0\).  Differentiating under the Laplace integral in \eqref{eq:reciprocal-gamma-kernel-representation} gives
\begin{equation}
\label{eq:reciprocal-gamma-alternating-integral}
(-1)^m\frac{d^m}{dx^m}H''(x)
=\int_0^\infty t^m e^{-xt}K(t)\,dt>0.
\end{equation}
This proves \eqref{eq:reciprocal-gamma-complete-monotonicity}.
\end{proof}

\section{External zeta inequality proofs}

\begin{theorem}[APP-0002: Affirmative solution of Nantomah zeta positivity problem]
\label{thm:nantomah-positivity}
For every \(n\in\N\),
\begin{equation}
\label{eq:nantomah-positivity}
(n+2)\zeta(n+1)\zeta(n+3)
-(n+1)\zeta(n+2)^2
-\zeta(n+1)\zeta(n+2)>0.
\end{equation}
\end{theorem}

\begin{proof}
Put \(s=n+1\ge2\) and \(Z_s=\zeta(s)\).  The expression in \eqref{eq:nantomah-positivity} becomes
\begin{equation}
\label{eq:nantomah-ks}
K_s=(s+1)Z_sZ_{s+2}-sZ_{s+1}^2-Z_sZ_{s+1}.
\end{equation}
Let \(X(m)=1/m\) under \(\rho_s\).  By \eqref{eq:zeta-law-moment},
\begin{equation}
\label{eq:nantomah-moments}
\E_s[X]=\frac{Z_{s+1}}{Z_s},
\qquad
\E_s[X^2]=\frac{Z_{s+2}}{Z_s}.
\end{equation}
Thus
\begin{equation}
\label{eq:nantomah-moment-form}
\frac{K_s}{Z_s^2}
=(s+1)\E_s[X^2]-s\E_s[X]^2-\E_s[X]
=(s+1)\Var_s(X)+\E_s[X]^2-\E_s[X].
\end{equation}

Write \(a=\zeta(s)-1\), \(b=\zeta(s+1)-1\), and \(c=\zeta(s+2)-1\).  Expanding gives
\begin{align}
\label{eq:nantomah-lq-split}
K_s&=L_s+Q_s,&
L_s&=sa+(s+1)c-(2s+1)b,&
Q_s&=(s+1)ac-sb^2-ab.
\end{align}
The linear part is termwise positive:
\begin{equation}
\label{eq:nantomah-linear-positive}
L_s=
\sum_{m=2}^{\infty}
m^{-s}\frac{(m-1)(s(m-1)-1)}{m^2}
\ge
2^{-s}\frac{s-1}{4}.
\end{equation}
Since \(b\le a/2\) and \(c\ge0\),
\begin{equation}
\label{eq:nantomah-quadratic-correction}
Q_s\ge -sb^2-ab\ge-\frac{s+2}{4}a^2.
\end{equation}
For \(s\ge4\),
\begin{equation}
\label{eq:nantomah-tail-bound}
a=\sum_{m=2}^{\infty}m^{-s}
\le2^{-s}+\int_2^\infty x^{-s}\,dx
=2^{-s}\left(1+\frac2{s-1}\right).
\end{equation}
Equations \eqref{eq:nantomah-linear-positive}, \eqref{eq:nantomah-quadratic-correction}, and \eqref{eq:nantomah-tail-bound} imply
\begin{equation}
\label{eq:nantomah-lower-bound}
K_s\ge
\frac14\left[
2^{-s}(s-1)
-(s+2)2^{-2s}\left(1+\frac2{s-1}\right)^2
\right].
\end{equation}
It is therefore enough to show
\begin{equation}
\label{eq:nantomah-sufficient-inequality}
2^s(s-1)>(s+2)\left(1+\frac2{s-1}\right)^2.
\end{equation}
Equivalently,
\begin{equation}
\label{eq:nantomah-f-function}
F(s)=\frac{2^s(s-1)^3}{(s+2)(s+1)^2}>1.
\end{equation}
Now \(F(4)=72/25>1\), and
\begin{equation}
\label{eq:nantomah-f-monotone}
\frac{d}{ds}\log F(s)
=\log2+\frac3{s-1}-\frac1{s+2}-\frac2{s+1}>0
\qquad(s\ge4).
\end{equation}
Thus \(K_s>0\) for \(s\ge4\).

For \(s=2\),
\begin{equation}
\label{eq:nantomah-k2}
K_2=3\zeta(2)\zeta(4)-2\zeta(3)^2-\zeta(2)\zeta(3).
\end{equation}
Using \(\zeta(3)<5/4\), \(\zeta(2)=\pi^2/6\), and \(\zeta(4)=\pi^4/90\),
\begin{equation}
\label{eq:nantomah-k2-bound}
K_2>
\frac{\pi^6}{180}-\frac{5\pi^2}{24}-\frac{25}{8}
>
\frac{10415360504209}{180000000000000}>0.
\end{equation}
For \(s=3\), the estimates \(\zeta(3)>251/216\), \(\zeta(5)>8051/7776\), and \(\zeta(4)<13/12\) yield
\begin{equation}
\label{eq:nantomah-k3-bound}
K_3>
\frac{251}{216}
\left(4\cdot\frac{8051}{7776}-\frac{13}{12}\right)
-3\left(\frac{13}{12}\right)^2
=\frac{13783}{419904}>0.
\end{equation}
Hence \(K_s>0\) for all integers \(s\ge2\), proving the theorem.  This answers Nantomah's question \cite{nantomah}.
\end{proof}

\begin{theorem}[APP-0001: Alzer-Kwong convexity and concavity pattern for reciprocal zeta]
\label{thm:alzer-kwong}
Let \(F(x)=1/\zeta(x)\).  For every integer \(n\ge1\),
\begin{equation}
\label{eq:alzer-kwong-convex}
F''(x)>0\quad\text{on }(-4n,-4n+2),
\end{equation}
and
\begin{equation}
\label{eq:alzer-kwong-concave}
F''(x)<0\quad\text{on }(-4n-2,-4n).
\end{equation}
\end{theorem}

\begin{proof}
Put \(x=-u\), where \(u>2\).  The functional equation gives
\begin{equation}
\label{eq:zeta-functional-equation-negative-axis}
\zeta(-u)
=-2^{-u}\pi^{-u-1}\sin\left(\frac{\pi u}{2}\right)\Gamma(u+1)\zeta(u+1),
\end{equation}
and hence
\begin{equation}
\label{eq:reciprocal-zeta-negative-axis}
\frac1{\zeta(-u)}
=-\frac{2^u\pi^{u+1}}
{\Gamma(u+1)\zeta(u+1)\sin(\pi u/2)}.
\end{equation}
Define the positive function
\begin{equation}
\label{eq:alzer-g-function}
G(u)=
\frac{2^u\pi^{u+1}}
{\Gamma(u+1)\zeta(u+1)|\sin(\pi u/2)|}.
\end{equation}
On \(u\in(4n-2,4n)\), \(\sin(\pi u/2)<0\), so \(F(-u)=G(u)\).  On \(u\in(4n,4n+2)\), \(\sin(\pi u/2)>0\), so \(F(-u)=-G(u)\).  Since second derivatives are unchanged by \(x=-u\), it remains to prove \(G''(u)>0\).

Let \(\ell(u)=\log G(u)\).  Then
\begin{equation}
\label{eq:alzer-log-g}
\ell(u)=u\log2+(u+1)\log\pi-\log\Gamma(u+1)-\log\zeta(u+1)
-\log\left|\sin\left(\frac{\pi u}{2}\right)\right|.
\end{equation}
Differentiating twice gives
\begin{equation}
\label{eq:alzer-log-g-second}
\ell''(u)
=-\psi'(u+1)-(\log\zeta)''(u+1)
+\frac{\pi^2}{4}\csc^2\left(\frac{\pi u}{2}\right).
\end{equation}
With \(s=u+1>3\),
\begin{equation}
\label{eq:trigamma-bound}
\psi'(s)=\sum_{k=0}^{\infty}\frac1{(s+k)^2}
<
\int_{s-1}^{\infty}\frac{dt}{t^2}
=\frac1{s-1}\le\frac12.
\end{equation}
Combining \eqref{eq:uniform-curvature-bound}, \eqref{eq:trigamma-bound}, and \(\csc^2(\pi u/2)\ge1\), we get
\begin{equation}
\label{eq:alzer-log-convexity}
\ell''(u)>
\frac{\pi^2}{4}-\frac12-\frac13
=\frac{\pi^2}{4}-\frac56>0.
\end{equation}
Since \(G>0\),
\begin{equation}
\label{eq:alzer-g-second-positive}
G''(u)=G(u)(\ell''(u)+\ell'(u)^2)>0.
\end{equation}
The two signs in \eqref{eq:alzer-kwong-convex} and \eqref{eq:alzer-kwong-concave} now follow from \(F(-u)=G(u)\) on \(u\in(4n-2,4n)\) and \(F(-u)=-G(u)\) on \(u\in(4n,4n+2)\).  This proves the Alzer--Kwong sign pattern \cite{alzer-kwong}.
\end{proof}

\begin{theorem}[APP-0003: Generalized Holder inequality for Gamma zeta]
\label{thm:sroysang-holder}
Let \(m\ge2\), \(r\ge1\), \(x_1,\ldots,x_m>1\), and \(p_1,\ldots,p_m>1\) satisfy
\begin{equation}
\label{eq:sroysang-holder-condition}
\sum_{i=1}^{m}\frac1{p_i}=\frac1r.
\end{equation}
Define
\begin{equation}
\label{eq:sroysang-t}
T=r\sum_{i=1}^{m}\frac{x_i}{p_i}.
\end{equation}
Then \(T>1\) and
\begin{equation}
\label{eq:sroysang-generalized-holder}
\zeta(T)\le
\frac{\prod_{i=1}^{m}(\Gamma(x_i)\zeta(x_i))^{r/p_i}}{\Gamma(T)}.
\end{equation}
\end{theorem}

\begin{proof}
Set
\begin{equation}
\label{eq:sroysang-lambda}
\lambda_i=\frac{r}{p_i}.
\end{equation}
Then \(\lambda_i>0\), \(\sum_i\lambda_i=1\), and \(T=\sum_i\lambda_i x_i>1\).  Define
\begin{equation}
\label{eq:sroysang-q}
q_i=\frac{p_i}{r}.
\end{equation}
Then \(\sum_i1/q_i=1\), and the assumptions imply \(q_i>1\) for all \(i\); otherwise one term \(1/p_i\) would already be at least \(1/r\), leaving no room for the remaining positive terms in \eqref{eq:sroysang-holder-condition}.

For \(t>0\), let
\begin{equation}
\label{eq:sroysang-fi}
f_i(t)=\left(\frac{t^{x_i-1}}{e^t-1}\right)^{1/p_i}.
\end{equation}
Then
\begin{equation}
\label{eq:sroysang-product}
\left(\prod_{i=1}^{m}f_i(t)\right)^r
=\frac{t^{r\sum_i(x_i-1)/p_i}}{e^t-1}
=\frac{t^{T-1}}{e^t-1}.
\end{equation}
Consequently,
\begin{equation}
\label{eq:sroysang-lr-norm}
\left\|\prod_{i=1}^{m}f_i\right\|_{L^r(0,\infty)}^r
=\int_0^\infty\frac{t^{T-1}}{e^t-1}\,dt
=\Gamma(T)\zeta(T).
\end{equation}
Generalized Holder gives
\begin{equation}
\label{eq:sroysang-holder-step}
\left\|\prod_{i=1}^{m}f_i\right\|_{L^r}
\le
\prod_{i=1}^{m}\|f_i\|_{L^{p_i}}.
\end{equation}
Moreover,
\begin{equation}
\label{eq:sroysang-pi-norm}
\|f_i\|_{L^{p_i}}^{p_i}
=\int_0^\infty\frac{t^{x_i-1}}{e^t-1}\,dt
=\Gamma(x_i)\zeta(x_i).
\end{equation}
Combining \eqref{eq:sroysang-lr-norm}, \eqref{eq:sroysang-holder-step}, and \eqref{eq:sroysang-pi-norm}, then raising to the power \(r\), gives
\begin{equation}
\label{eq:sroysang-predivide}
\Gamma(T)\zeta(T)
\le
\prod_{i=1}^{m}(\Gamma(x_i)\zeta(x_i))^{r/p_i}.
\end{equation}
Division by \(\Gamma(T)\) proves \eqref{eq:sroysang-generalized-holder}.  This solves the generalized Holder problem in Sroysang's notation, where his \(\xi(s)\) is \(\zeta(s)\) for \(s>1\) \cite{sroysang}.
\end{proof}

\begin{remark}[Free-energy interpretation of the Sroysang inequality]
The same result follows from log-convexity of \(M(s)=\Gamma(s)\zeta(s)\).  With \(\lambda_i=r/p_i\) and \(T=\sum_i\lambda_ix_i\), Jensen's inequality for \(\log M\) gives
\begin{equation}
\label{eq:sroysang-jensen}
M(T)\le\prod_{i=1}^{m}M(x_i)^{\lambda_i}
=\prod_{i=1}^{m}(\Gamma(x_i)\zeta(x_i))^{r/p_i},
\end{equation}
which is equivalent to \eqref{eq:sroysang-generalized-holder}.
\end{remark}

\section{Conclusion}

The zeta-law normalization gives a single language for several identities and inequalities involving \(\zeta\).  At the structural level, the paper proves the successor entropy formula \eqref{eq:successor-entropy-series} and its modular resolution \eqref{eq:successor-entropy-modular-resolution}, the Euler-score decomposition \eqref{eq:euler-score-decomposition}, the Mellin-Planck variance identity \eqref{eq:mellin-planck-log-convexity}, the reciprocal-Gamma complete-monotonicity theorem \eqref{thm:reciprocal-gamma-complete-monotonicity}, and the tail partition formulas \eqref{eq:tail-s7-floor} and \eqref{eq:tail-s8-floor}.  These results are not only formal translations: they supply direct estimates for the six labelled applications in Section~\ref{sec:applications}.  APP-0001 follows from the functional equation and a positive-axis curvature estimate; APP-0002 follows from a moment split and tail bounds; APP-0003 follows from the Planck-kernel integral representation; APP-0004 follows from a positive Laplace kernel; APP-0005 follows from a rational telescoping enclosure; and APP-0006 follows from adjacent corrected asymptotic truncations.  The remaining useful directions are to understand whether the reciprocal-Weierstrass kernel method extends to products of higher polygamma functions and to sharpen the modular entropy side, especially the prime-modulus functionals in \eqref{eq:prime-functional-limit}, where additive successor entropy is encoded by finite Dirichlet \(L\)-value data.

\begin{thebibliography}{10}

\bibitem{nantomah}
K. Nantomah,
\emph{Open Problem on Riemann Zeta Function},
ResearchGate problem note, October 2024.
\url{https://www.researchgate.net/publication/384676538_Open_Problem_on_Riemann_Zeta_Function}.

\bibitem{sroysang}
B. Sroysang,
\emph{Two Inequalities for the Riemann Zeta Functions},
Mathematica Aeterna 3, no. 1 (2013), 21--24.
\url{https://arastirmax.com/en/system/files/dergiler/135290/makaleler/3/1/arastirmax-two-inequalities-riemann-zeta-functions.pdf}.

\bibitem{alzer-kwong}
H. Alzer and M. K. Kwong,
\emph{On the concavity and convexity of \(1/\zeta\)},
International Journal of Number Theory 21, no. 8 (2025), 1825--1835.
\url{https://doi.org/10.1142/S1793042125500897}.

\bibitem{qi-lim-nantomah-polygamma}
F. Qi, D. Lim, and K. Nantomah,
\emph{Monotonicity and positivity of several functions involving ratios and products of polygamma functions},
Journal of Inequalities and Applications 2025, article 5.
\url{https://doi.org/10.1186/s13660-024-03245-8}.

\bibitem{kim-song-tail}
D. Kim and K. Song,
\emph{The inverses of tails of the Riemann zeta function},
Journal of Inequalities and Applications 2018, article 157.
\url{https://doi.org/10.1186/s13660-018-1743-6}.

\bibitem{pan-wu-tail}
Z. Pan and Z. Wu,
\emph{The inverse of tails of Riemann zeta function, Hurwitz zeta function and Dirichlet L-function},
AIMS Mathematics 9, no. 6 (2024), 16564--16585.
\url{https://doi.org/10.3934/math.2024803}.

\bibitem{apostol}
T. M. Apostol,
\emph{Introduction to Analytic Number Theory},
Undergraduate Texts in Mathematics, Springer, 1976.
\url{https://doi.org/10.1007/978-1-4757-5579-4}.

\bibitem{davenport}
H. Davenport,
\emph{Multiplicative Number Theory},
3rd ed., revised by H. L. Montgomery, Graduate Texts in Mathematics, vol. 74, Springer, 2000.
\url{https://doi.org/10.1007/978-1-4757-5927-3}.

\bibitem{titchmarsh}
E. C. Titchmarsh,
\emph{The Theory of the Riemann Zeta-function},
2nd ed., revised by D. R. Heath-Brown, Oxford University Press, 1986.
\url{https://global.oup.com/academic/product/the-theory-of-the-riemann-zeta-function-9780198533696}.

\bibitem{cover-thomas}
T. M. Cover and J. A. Thomas,
\emph{Elements of Information Theory},
2nd ed., Wiley-Interscience, 2006.
\url{https://doi.org/10.1002/047174882X}.

\end{thebibliography}

\end{document}
